\documentclass[12pt,a4paper]{article}

\usepackage[utf8]{inputenc}
\usepackage[T1]{fontenc}
\usepackage[brazil]{babel}
\usepackage{geometry}
\usepackage{listings}
\usepackage{xcolor}
\usepackage{longtable}
\usepackage{fancyvrb}

\geometry{margin=2.5cm}

\title{SQL (Structured Query Language) e DB (Database)}
\author{}
\date{}

\begin{document}

\maketitle

\begin{itemize}
    \item Antes de definir o que é SQL, precisamos definir o que é um banco de dados.
    
    \item Um banco de dados é um conjunto de informações organizadas que estão armazenadas em um sistema (manual ou eletrônico).
    
    \item Um dado é a menor unidade de informação, o que compõe uma tabela. Um conjunto de tabelas forma um banco de dados.
\end{itemize}

A definição pode ser observada abaixo:

\begin{verbatim}
====================================================================
DADO
====================================================================

Um dado é uma informação única e isolada.

Exemplos:
--------------------------------------------------
"Jill"
22
"Brasil"
true
95.7
--------------------------------------------------

Visualmente:

+------------------+
| DADO             |
+------------------+
| Nome: Jill       |
+------------------+

Outro exemplo:

+------------------+
| DADO             |
+------------------+
| Idade: 22        |
+------------------+


====================================================================
TABELA
====================================================================

Uma tabela organiza vários dados relacionados
em linhas e colunas.

Exemplo:
Tabela de usuários.

+----+----------+-------+----------------------+
| ID | NOME     | IDADE | EMAIL                |
+----+----------+-------+----------------------+
| 1  | Ana      | 20    | ana@email.com        |
| 2  | Carlos   | 31    | carlos@email.com     |
| 3  | Jill     | 22    | jill@email.com       |
+----+----------+-------+----------------------+

Estrutura:

COLUNAS
   ↓

+----+----------+-------+
| ID | NOME     | IDADE |
+----+----------+-------+
| 1  | Ana      | 20    | ← LINHA / REGISTRO
| 2  | Carlos   | 31    |
+----+----------+-------+

Cada célula:
-------
| Ana |
-------

é um DADO.


====================================================================
BANCO DE DADOS
====================================================================

Um banco de dados é um conjunto organizado
de tabelas relacionadas.

Visualização:

                 BANCO DE DADOS
          "Sistema da Empresa"

     +-----------------------------+
     | TABELA: usuarios            |
     +-----------------------------+
     | id | nome | idade | email  |
     +-----------------------------+

     +-----------------------------+
     | TABELA: produtos            |
     +-----------------------------+
     | id | nome | preco           |
     +-----------------------------+

     +-----------------------------+
     | TABELA: pedidos             |
     +-----------------------------+
     | id | usuario_id | produto_id|
     +-----------------------------+


====================================================================
RELAÇÃO ENTRE AS PARTES
====================================================================

DADO
 ↓
Menor unidade de informação

Ex:
----------
| "Jill" |
----------


VÁRIOS DADOS
 ↓
formam uma LINHA

+----+------+-------+
| 1  | Jill | 22    |
+----+------+-------+


VÁRIAS LINHAS
 ↓
formam uma TABELA

+----+------+-------+
| ID | Nome | Idade | -> Atributo, coluna ou campo.
+----+------+-------+
| 1  | Jill | 22    | -> Tupla, linha ou registros.
| 2  | Ana  | 20    |
+----+------+-------+


VÁRIAS TABELAS
 ↓
formam um BANCO DE DADOS

+----------------------+
| Banco de Dados       |
+----------------------+
| usuarios             |
| produtos             |
| pedidos              |
+----------------------+


====================================================================
EXEMPLO COMPLETO
====================================================================

BANCO DE DADOS: Escola
│
├── TABELA alunos
│      ├── dado: "João"
│      ├── dado: 18
│      └── dado: "3A"
│
├── TABELA professores
│      ├── dado: "Maria"
│      └── dado: "Matemática"
│
└── TABELA notas
       ├── dado: 9.5
       └── dado: 7.0


====================================================================
RESUMO FINAL
====================================================================

+----------------+----------------------------------+
| CONCEITO       | DEFINIÇÃO                        |
+----------------+----------------------------------+
| DADO           | Informação isolada               |
| TABELA         | Conjunto organizado de dados     |
| BANCO DE DADOS | Conjunto organizado de tabelas   |
+----------------+----------------------------------+
====================================================================
\end{verbatim}

\section*{Sistema Gerenciador de Banco de Dados}

Um SGBD é uma coleção de dados inter-relacionados e um conjunto de softwares que acessam esses dados. O SGBD proporciona um ambiente para recuperar e armazenar informações no banco de dados. Esse sistema contém 4 componentes básicos, sendo:

\begin{itemize}
    \item \textbf{Hardware}: São os meios em que os dados são armazenados.
    
    \item \textbf{Dados}: Os dados armazenados são repartidos e podem ser acessados por usuários diferentes.
    
    \item \textbf{Software}: Sistema que se encontra entre o usuário e o hardware, que serve unicamente para gerenciar o banco de dados.
    
    \item \textbf{Usuários}: Existem três classes de usuários, sendo o primeiro responsável por escrever programas de aplicação que usam o banco de dados; a segunda classe é o usuário final que acessa o banco de dados através de um terminal que faz o uso de uma linguagem para fazer o acesso, essa linguagem pode ser o SQL. E a última classe é o administrador de banco de dados (DBA). Ele analisa quais informações podem ser mantidas, identifica as entidades de interesse das empresas e é responsável pelo desempenho do sistema.
\end{itemize}

\section*{Modelo Relacional}

Nesse modelo de banco de dados as relações acontecem através de tabelas. Uma tabela possui linhas e colunas e todas as tabelas devem possuir um nome e um conjunto de atributos/campos. As colunas são os atributos da tabela e devem possuir um nome e o tipo de dado que será armazenado na coluna. Cada conjunto de atributos forma uma linha/registro, que também pode ser chamado de tupla. Nesse modelo os dados devem estar organizados e separados de acordo com suas relações.

\begin{verbatim}
====================================================================
BANCO DE DADOS RELACIONAL (SQL)
====================================================================

Os dados ficam organizados em TABELAS
relacionadas entre si.

Usa:
- linhas
- colunas
- chaves
- relacionamentos

Exemplos de sistemas de gerenciamento:
--------------------------------------------------
MySQL
PostgreSQL
Oracle
SQL Server
SQLite
--------------------------------------------------
\end{verbatim}

\section*{Modelo Não Relacional}

Modelo mais atual que permite o acesso aos dados fora das estruturas encontradas no banco de dados relacionais. O NoSQL pode acessar dados do banco relacional. Essa estrutura, ao invés de armazenar os dados de forma tabular, os armazena em um documento em forma de JSON.

\begin{verbatim}
====================================================================
BANCO DE DADOS NÃO RELACIONAL (NoSQL)
====================================================================

Os dados NÃO precisam ficar em tabelas rígidas.

Os formatos podem variar.

Pode usar:
- documentos
- chave/valor
- grafos
- colunas

Exemplos de sistemas de gerenciamento:
--------------------------------------------------
MongoDB
Redis
Cassandra
Neo4j
Firebase
--------------------------------------------------
\end{verbatim}

\section*{Terminal}

Nos softwares de gerenciamento, toda a interação será feita através de um terminal (CLI). No terminal é usado texto puro para enviar instruções aos SGBDs. É através do terminal que os DBAs realizam consultas, atualizações, importação de dados, e ainda permite uma conexão segura com os bancos de dados via SSH. É importante salientar que o terminal é a parte visual onde são digitados os comandos; o shell é o programa que roda no terminal em que digitamos, como no meu caso, o bash é o modelo de shell que eu uso, padrão na maioria das distribuições Linux.

\section*{A Linguagem SQL}

Embora seja considerada uma linguagem de consulta, ela também possui recursos para definição de estruturas de dados, para modificar dados no banco e recursos para especificar restrições de segurança e identidade. SQL pode ser definida como uma linguagem de uso especial, sendo usada somente para gerenciar banco de dados relacional. No ano de 1986 a ANSI definiu a linguagem SQL como padrão e documentou centenas de especificações detalhadas sobre como uma linguagem de banco de dados deve funcionar.

A linguagem possui diversas partes, sendo as principais:

\begin{itemize}
    \item \textbf{Data Definition Language (DDL)}: Fornece comandos para definir e modificar os esquemas de relação, remoção e criação de índices. Comandos principais: CREATE, ALTER, DROP.
    
    \item \textbf{Data Manipulation Language (DML)}: Modelo de busca baseada em álgebra e cálculo relacional, e inclui comandos para inserir, remover e modificar informações. Comandos principais: SELECT, INSERT, UPDATE, DELETE.
    
    \item \textbf{Data Control Language (DCL)}: Conjunto de comandos que realiza o cadastro de usuários e determina o nível de privilégio para objetos no banco de dados. Comandos principais: GRANT, REVOKE.
    
    \item \textbf{Transactions Control}: Possui comandos que especificam o início e o fim das transações. Algumas implementações permitem o trancamento explícito dos dados para controle de concorrência. Comandos: COMMIT, ROLLBACK, SAVEPOINT.
\end{itemize}

\section*{Comandos SQL - Tabela Resumida}

\begin{verbatim}
===========================================================
               COMANDOS SQL - TABELA RESUMIDA
===========================================================

-----------------------------------------------------------
CRIAÇÃO DE BANCO DE DADOS
-----------------------------------------------------------

CRIAR BANCO:
CREATE DATABASE nome_banco;

USAR BANCO:
USE nome_banco;

EXCLUIR BANCO:
DROP DATABASE nome_banco;

-----------------------------------------------------------
CRIAÇÃO DE TABELAS
-----------------------------------------------------------

CRIAR TABELA:
CREATE TABLE usuarios (
    id INT PRIMARY KEY,
    nome VARCHAR(100),
    idade INT
);

EXCLUIR TABELA:
DROP TABLE usuarios;

RENOMEAR TABELA:
RENAME TABLE usuarios TO clientes;

-----------------------------------------------------------
ALTERAÇÃO DE TABELAS
-----------------------------------------------------------

ADICIONAR COLUNA:
ALTER TABLE usuarios
ADD email VARCHAR(100);

MODIFICAR COLUNA:
ALTER TABLE usuarios
MODIFY nome VARCHAR(200);

RENOMEAR COLUNA:
ALTER TABLE usuarios
RENAME COLUMN nome TO nome_completo;

REMOVER COLUNA:
ALTER TABLE usuarios
DROP COLUMN email;

-----------------------------------------------------------
INSERÇÃO DE DADOS
-----------------------------------------------------------

INSERIR DADOS:
INSERT INTO usuarios (id, nome, idade)
VALUES (1, 'Jill', 20);

INSERIR MÚLTIPLAS LINHAS:
INSERT INTO usuarios
VALUES
(2, 'Ana', 19),
(3, 'Carlos', 25);

-----------------------------------------------------------
CONSULTA DE DADOS
-----------------------------------------------------------

SELECIONAR TODOS:
SELECT * FROM usuarios;

SELECIONAR COLUNAS:
SELECT nome, idade FROM usuarios;

FILTRAR DADOS:
SELECT * FROM usuarios
WHERE idade > 18;

ORDENAR:
SELECT * FROM usuarios
ORDER BY nome ASC;

LIMITAR RESULTADOS:
SELECT * FROM usuarios
LIMIT 5;

-----------------------------------------------------------
OPERADORES
-----------------------------------------------------------

IGUAL:
WHERE idade = 20

DIFERENTE:
WHERE idade != 20

MAIOR:
WHERE idade > 20

MENOR:
WHERE idade < 20

ENTRE:
WHERE idade BETWEEN 18 AND 30

LISTA:
WHERE cidade IN ('SP', 'RJ')

PESQUISA:
WHERE nome LIKE '%an%'

-----------------------------------------------------------
ATUALIZAÇÃO DE DADOS
-----------------------------------------------------------

ATUALIZAR:
UPDATE usuarios
SET idade = 21
WHERE id = 1;

-----------------------------------------------------------
REMOÇÃO DE DADOS
-----------------------------------------------------------

REMOVER LINHA:
DELETE FROM usuarios
WHERE id = 1;

REMOVER TODOS OS DADOS:
TRUNCATE TABLE usuarios;

-----------------------------------------------------------
CHAVES
-----------------------------------------------------------

PRIMARY KEY:
id INT PRIMARY KEY

FOREIGN KEY:
CREATE TABLE pedidos (
    id INT PRIMARY KEY,
    usuario_id INT,
    FOREIGN KEY (usuario_id)
    REFERENCES usuarios(id)
);

-----------------------------------------------------------
JOINS
-----------------------------------------------------------

INNER JOIN:
SELECT usuarios.nome, pedidos.id
FROM usuarios
INNER JOIN pedidos
ON usuarios.id = pedidos.usuario_id;

LEFT JOIN:
SELECT *
FROM usuarios
LEFT JOIN pedidos
ON usuarios.id = pedidos.usuario_id;

RIGHT JOIN:
SELECT *
FROM usuarios
RIGHT JOIN pedidos
ON usuarios.id = pedidos.usuario_id;

-----------------------------------------------------------
FUNÇÕES
-----------------------------------------------------------

CONTAR:
SELECT COUNT(*) FROM usuarios;

MÉDIA:
SELECT AVG(idade) FROM usuarios;

MÁXIMO:
SELECT MAX(idade) FROM usuarios;

MÍNIMO:
SELECT MIN(idade) FROM usuarios;

SOMA:
SELECT SUM(idade) FROM usuarios;

-----------------------------------------------------------
AGRUPAMENTO
-----------------------------------------------------------

AGRUPAR:
SELECT cidade, COUNT(*)
FROM usuarios
GROUP BY cidade;

FILTRAR GRUPOS:
SELECT cidade, COUNT(*)
FROM usuarios
GROUP BY cidade
HAVING COUNT(*) > 5;

-----------------------------------------------------------
ÍNDICES
-----------------------------------------------------------

CRIAR ÍNDICE:
CREATE INDEX idx_nome
ON usuarios(nome);

REMOVER ÍNDICE:
DROP INDEX idx_nome;

-----------------------------------------------------------
USUÁRIOS E PERMISSÕES
-----------------------------------------------------------

CRIAR USUÁRIO:
CREATE USER 'jill' IDENTIFIED BY '123';

DAR PERMISSÃO:
GRANT ALL PRIVILEGES
ON banco.*
TO 'jill';

REMOVER PERMISSÃO:
REVOKE ALL PRIVILEGES
ON banco.*
FROM 'jill';

-----------------------------------------------------------
COMANDOS IMPORTANTES
-----------------------------------------------------------

MOSTRAR BANCOS:
SHOW DATABASES;

MOSTRAR TABELAS:
SHOW TABLES;

DESCREVER TABELA:
DESCRIBE usuarios;

-----------------------------------------------------------
TIPOS DE DADOS
-----------------------------------------------------------

INT            -> Número inteiro
FLOAT          -> Número decimal
VARCHAR(n)     -> Texto variável
CHAR(n)        -> Texto fixo
TEXT           -> Texto longo
DATE           -> Data
TIME           -> Hora
DATETIME       -> Data e hora
BOOLEAN        -> Verdadeiro/Falso

-----------------------------------------------------------
BOAS PRÁTICAS
-----------------------------------------------------------

- Sempre usar WHERE em UPDATE e DELETE
- Utilizar PRIMARY KEY
- Criar índices em colunas muito pesquisadas
- Nomear tabelas claramente
- Evitar SELECT * em sistemas grandes
- Fazer backup regularmente

===========================================================
\end{verbatim}

\textbf{OBS: Este documento é apenas para fins de documentar o aprendizado. Não tem finalidades acadêmicas ou científicas.}

\section*{Fontes}

\begin{itemize}
    \item ZHAO, Alice. \textit{SQL: Um guia prático para o uso de SQL}. NovaTec/SP. 2022.
    
    \item FERNEDA, Edberto. \textit{Introdução à Linguagem SQL}. USP/SP.
    
    \item \texttt{https://www.ibm.com/br-pt/think/topics/database}
\end{itemize}

\end{document}
